\subsection{The hammock comparison and the scope of the remaining inclusion problem}

\begin{result}{B.1}{Proved}{The two diagrams and the actual inclusion}
For a type $\tau=\varepsilon_1\cdots\varepsilon_n$, write $I_\tau$ for the free category on vertices $0,\ldots,n$ and the directed edges prescribed by its signs. A type map is a nondecreasing endpoint-preserving functor. Equivalently, its $i$th edge is sent to a string of $l_i\geq0$ edges of the same sign; the target word is $\varepsilon_1^{l_1}\cdots\varepsilon_n^{l_n}$. Let $\Sigma$ be a wide subcategory of a small category $M$.

Let $F_\tau(a,b)$ be the full category of functors $I_\tau\to M$ with endpoints $a,b$ and backward edges in $\Sigma$, with transformations equal to identities at both endpoints. Let $F_\tau^\Sigma(a,b)$ have the same objects and only transformations whose components all lie in $\Sigma$. These are strict contravariant diagrams on $\mathrm{Ty}$. Their fibred Grothendieck constructions are respectively
$\operatorname{Ch}_\Sigma(M;a,b)$ and
$\operatorname{Ch}_\Sigma^\Sigma(M;a,b)$.
The inclusion of diagrams induces the specified identity-on-objects functor
\[
 J:\operatorname{Ch}_\Sigma^\Sigma(M;a,b)
       \longrightarrow\operatorname{Ch}_\Sigma(M;a,b).
\]
Both constructions and $J$ are strictly compatible with concatenation and empty-path units.
\end{result}
\begin{proof}
An arrow $(\tau,c)\to(\sigma,d)$ is a pair $(u,\eta)$ with
$u:I_\tau\to I_\sigma$ a type map and
$\eta:c\Rightarrow d u$ an endpoint-fixed natural transformation. In the restricted category, every component of $\eta$ must lie in $\Sigma$. Its composite with $(v,\theta):(\sigma,d)\to(\rho,e)$ is
$(vu,(\theta u)\eta)$.
Composition closure of $\Sigma$ proves that these arrows form the restricted category. Precomposition sends a backward edge to an identity or to a composite of backward edges, so both diagrams preserve their stated object conditions. It also preserves the vertical restriction.

Concatenation glues the two endpoints literally and concatenates the two type maps and the two component lists. The common component is an identity. The displayed composition formula commutes with this operation. Concatenation of finite ordered lists is strictly associative, and the empty list is a strict unit. This verifies the claims without identifying an equivalence of classifying spaces with an equivalence of categories.
\end{proof}

\begin{result}{B.2}{Known}{A strict and composition-compatible hammock reduction}
For every small relative category with a wide subcategory $\Sigma$, there is a specified simplicial map
\[
 \rho_{a,b}:N\operatorname{Ch}_\Sigma^\Sigma(M;a,b)
       \longrightarrow L^H(M,\Sigma)(a,b)
\]
which is a weak homotopy equivalence. It is natural in relative functors and forms, as $a,b$ vary, a simplicial functor that is the identity on objects and strictly respects concatenation and units. No two-out-of-three or saturation assumption is needed.
\end{result}
\begin{proof}
The fixed-shape restricted category is exactly the category of width-zero and width-one hammocks of Dwyer and Kan, Definition~5.1. Their indexing category $\mathbf {II}$ is $\mathrm{Ty}^{\mathrm{op}}$, with the word-reading convention reversed if one writes composites in the opposite order. Proposition~5.5 identifies the hammock mapping simplicial set with
\[
 \operatorname*{colim}_{\tau\in\mathrm{Ty}^{\mathrm{op}}}
             N F_\tau^\Sigma(a,b),                         \tag{B.2.1}
\]
where the colimit is in simplicial sets, not in categories. Their Proposition~2.2 identifies this localization with simplicial localization \citep{DwyerKan1980Calculating}.

We specify the map from the manuscript's fibred convention. An $n$-simplex of its nerve is a string of path morphisms with types
$\tau_0\to\tau_1\to\cdots\to\tau_n$. Pull all $n+1$ paths back to $\tau_0$ and use the given transformations to form a width-$n$ restricted hammock. Send it to (B.2.1). All face operations except the zeroth simply compose or omit a row. The zeroth face changes the chosen common type from $\tau_0$ to $\tau_1$; the two resulting simplices agree in (B.2.1) by the type-change relation. Degeneracies insert identity rows. Thus this is an actual simplicial map $\rho$, not merely a map after localization. Pullback to the first type and reduction commute with concatenating two strings of path morphisms, proving the strict composition assertion.

For clarity, we verify the homotopy-colimit issue in this exact comparison. The type category is Reedy with degree the number of edges, direct maps injective on vertices, and inverse maps surjective on vertices. The factorization is through the image on vertices. In $\mathrm{Ty}^{\mathrm{op}}$, latching maps for $\tau\mapsto NF_\tau^\Sigma$ are the inclusions of the union of the subsets in which one or more specified horizontal edges are identities. Intersections correspond to collapsing the union of those edge sets. Hence these are monomorphisms and the diagram is Reedy cofibrant.

The Reedy category $\mathrm{Ty}^{\mathrm{op}}$ has fibrant constants. Indeed, a proper injective endpoint-preserving map into $I_\tau$ has to contain every turning vertex and both endpoints. If $\tau$ is not reduced, the inclusion of this reduced subinterval is initial among such maps; if it is reduced, there are no proper such maps. The relevant inverse latching category is therefore empty or connected. Consequently ordinary colimit computes the derived colimit of this diagram. This is the argument of Arakawa and Cnossen, Proposition~4.8 and Lemma~4.16; their Theorem~4.7 gives the simplicial-category comparison, and Definition~1.7, Theorem~1.9, and Corollary~1.11 identify the resulting anima with the mapping space of localization \citep{ArakawaCnossenPATHS}. Their Remark~4.18 explicitly records the Grothendieck-construction variant.

Finally the natural Grothendieck homotopy-colimit comparison identifies
$B\operatorname{Ch}_\Sigma^\Sigma$ with the homotopy colimit of these fibre nerves. The map to (B.2.1) is its canonical colimit cocone, as is also seen from the first-type formula above. The Reedy calculation proves that this specified map is a weak equivalence. For the natural Grothendieck comparison one may use Chach\'olski and Scherer, Theorem~26.8 and Corollary~26.9 \citep{ChacholskiSchererPATHS}. The Reedy proof and Dwyer--Kan's original definition apply to any wide subcategory; the convention that a relative $\infty$-category contains all equivalences in the more recent paper is not needed to add an unstated hypothesis here.
\end{proof}

\begin{result}{B.3}{Refuted}{Two-out-of-three does not make $J$ an isomorphism}
Even for saturated $\Sigma$, $J$ need not be an isomorphism of categories. It already fails for $M=[1]=\{0\to1\}$ and $\Sigma$ consisting of identities.
\end{result}
\begin{proof}
Take endpoints $0,1$, type $(+,+)$, and paths
$p=(0,0,1)$ and $q=(0,1,1)$. There is a full path transformation
$p\Rightarrow q$ whose components are
$(1_0,0\to1,1_1)$. Its squares commute. This transformation is not in the restricted category. There are two objects and three arrows in $M$; the only composites involve identities. The identity class is wide, satisfies two-out-of-three, and is saturated because localization changes nothing and the nonidentity arrow is not invertible. The finite path-category enumerator includes this category and retains the distinct arrows determined by type maps. This example refutes only the claimed categorical isomorphism: B.4 proves that its $J$ is a weak equivalence.
\end{proof}

\begin{result}{B.4}{Proved}{Denominators that are already isomorphisms}
Suppose every arrow of $\Sigma$ is an isomorphism in $M$; $\Sigma$ need not contain every isomorphism. Then both path categories have homotopy-discrete nerves, with components canonically $\operatorname{Hom}_M(a,b)$, and $J$ is a weak equivalence. The full all-type category is already isomorphic to this discrete category in $\mathrm{Cat}\,\mathrm{hot}$. In particular Lead~2 holds for $\Sigma=\operatorname{Iso}(M)$. For $\Sigma=\operatorname{Mor}(M)$, $J$ is the identity and B.2 applies directly.
\end{result}
\begin{proof}
Invert the backward arrows in $M$ and compose. This defines a functor
$e:\operatorname{Ch}_\Sigma(M;a,b)\to\operatorname{Hom}_M(a,b)_{\mathrm{disc}}$.
Naturality of a ladder, with identity endpoint components, implies equality of its two effects. Precomposition by a type map leaves the effect unchanged. Let $i$ insert an arrow as a one-edge forward path. Then $ei=1$.

There is a uniform contraction, stronger than contractibility after realization. For a path $c$ of type $\tau$, let $\operatorname{pad}(c)$ append a forward identity edge at $b$, so its type is $\tau*(+)$. On arrows use the old type map concatenated with the identity of $(+)$, and append the identity endpoint component. Collapsing the appended edge gives a natural transformation $\operatorname{pad}\Rightarrow1$.

Define $T(c)$, also of type $\tau*(+)$, to be $a$ at every old vertex, $b$ at the new last vertex, the identity on every old edge, and $e(c)$ on the appended forward edge. On a path morphism $(u,\eta)$ use $u*1_{(+)}$ and identity components; the two effects are equal, so this is a functor. The accumulated prefix effects $a\to c_j$, calculated by inverting the backward arrows, give an endpoint-fixed transformation $T(c)\Rightarrow\operatorname{pad}(c)$. The old last component is $e(c):a\to b$ and the new last component is $1_b$, so the final square also commutes. A path morphism preserves every accumulated prefix after reparametrization; this follows by induction across its commuting squares, using inverse backward arrows. Thus these transformations are natural in $c$. Collapsing the entire old word to the source of the appended edge gives a natural transformation $T\Rightarrow ie$ with identity components. We obtain the finite zigzag
\[
       1\ \Longleftarrow\ \operatorname{pad}
          \ \Longleftarrow\ T\ \Longrightarrow\ ie.
\]
It proves that $e$ and $i$ are inverse in $\mathrm{Cat}\,\mathrm{hot}$, not merely that $Be$ is a weak equivalence. Empty paths and empty Hom-sets satisfy the same formulas, with the unique functor on an empty category in the latter case.

Inverting arrows that are already equivalences does not change the $\infty$-category $NM$. B.2 therefore identifies the restricted nerve with the same discrete mapping space. Its effect map is the same canonical map on components. Since $J$ commutes with the two effect maps, two-out-of-three for weak homotopy equivalences proves the claim. The two-out-of-three principle here concerns maps of spaces, not an additional requirement on $\Sigma$.

The parenthetical explanation in Lead~2 cannot mean that the one-edge path is an initial object of the total all-type category: there need not be a type map $(+)\to\tau$. The displayed uniform contraction supplies the required correction and establishes the stronger statement in $\mathrm{Cat}\,\mathrm{hot}$.
\end{proof}

\begin{result}{B.5}{Known}{What a calculus of fractions identifies directly}
If $(M,\Sigma)$ admits the homotopy calculus of two-sided fractions of Dwyer--Kan Definition~6.1, the nerve of the restricted three-arrow category of diagrams
\[
       a\xleftarrow{\Sigma}x\longrightarrow y
             \xleftarrow{\Sigma}b
\]
computes $L^H(M,\Sigma)(a,b)$. If $\Sigma$ has two-out-of-three, the full and restricted categories of these three-arrow diagrams are equal. Corresponding statements hold for the two-arrow fraction categories under the left or right calculus, with the variance of Dwyer--Kan's notation retained. These statements alone do not identify the inclusion of the three-arrow category into the full all-type path category.
\end{result}
\begin{proof}
The first assertion is exactly Dwyer and Kan, Proposition~6.2(i), with the one-sided versions in (ii)--(iii) \citep{DwyerKan1980Calculating}. A vertical component at $x$ lies in $\Sigma$ because its composite with the upper or lower arrow to $a$ is the other arrow to $a$. The same argument, on the arrow from $b$, handles the component at $y$. This proves equality of the two three-arrow categories when two-out-of-three holds. It also proves the one-sided equalities with the appropriate single interior component.

Model categories give the standard homotopy-function-complex interpretation of the hammock side: Dwyer and Kan, \emph{Function complexes}, Propositions~4.4--4.5 and Corollary~4.7, identify it with resolution complexes and, for a cofibrant source and fibrant target in a simplicial model category, the simplicial mapping object \citep{DwyerKan1980Function}. For categories of fibrant objects the fraction statement is Arakawa--Cnossen, Theorem~3.26, which assumes their stated Brown-type axioms and requires no functorial factorization \citep{ArakawaCnossenPATHS}.

All these comparisons concern the restricted hammock construction or the specified short fraction category. The implication to an arbitrary-vertical all-type category would require proving that its additional inclusion is a weak equivalence. Neither endpoint forcing in the three-arrow shape nor Proposition~6.2 proves that additional assertion. We retain this distinction in B.8.
\end{proof}

\begin{result}{B.6}{Proved}{The full category has a canonical map to the localization mapping space}
For $\mathcal C=L(M,\Sigma)$ and $a,b\in M$, there is a natural map of spaces
\[
 Q_{a,b}:B\operatorname{Ch}_\Sigma(M;a,b)
           \longrightarrow\operatorname{Map}_{\mathcal C}(a,b)
\]
compatible with concatenation. Its composite with $BJ$ is the canonical hammock comparison of B.2. Consequently $BJ$ has a left inverse in $\mathrm{Hot}$ and induces injections on homotopy groups at every source basepoint. Both categories have the same component set
$\operatorname{Hom}_{M[\Sigma^{-1}]}(a,b)$, and $J$ induces that identity bijection.
\end{result}
\begin{proof}
For each type $\tau$, let $p(\tau)$ be the number of forward edges and let
$q_\tau:I_\tau\to[p(\tau)]$ count the forward edges preceding a vertex. It is the localization of the interval at its backward edges. This elementary interval localization is also specified in Arakawa--Cnossen, Notation~1.2 \citep{ArakawaCnossenPATHS}. Precomposition therefore gives an equivalence from $\operatorname{Fun}([p(\tau)],\mathcal C)$ to the full subcategory of $\operatorname{Fun}(I_\tau,\mathcal C)$ sending backward edges to equivalences. It commutes with the two endpoint evaluations. On their $\infty$-categorical fibres at $(a,b)$, composing along $0\to p(\tau)$ gives a functor to the space $\operatorname{Map}_{\mathcal C}(a,b)$.

We make the simultaneous coherence explicit. The ordinals $[p(\tau)]$ form a functor on $\mathrm{Ty}$: a type map sends each forward edge to its same-sign forward string, and the resulting monotone map of ordinals satisfies $q_\sigma u=\overline u q_\tau$. The endpoint arrows $[1]\to[p(\tau)]$ commute with all these maps. Therefore the precomposition equivalences above assemble to an equivalence of diagrams on $\mathrm{Ty}^{\mathrm{op}}$ with values in $\mathrm{Cat}_\infty$, after taking the specified endpoint fibres. We invert this \emph{equivalence of diagrams}, rather than choose inverses separately at each type. Inversion with its unit and counit is contractible in the functor $\infty$-category. The nerve of $F_\tau(a,b)$ maps naturally to the relative-diagram fibre, because its endpoint components are identities. Composing this map of diagrams with the inverse diagram equivalence and the endpoint-arrow composition map gives a coherent map to the constant diagram $\operatorname{Map}_{\mathcal C}(a,b)$. Groupoid completion, pointwise, produces a cocone on $BF_\tau$. Taking its colimit and the natural Grothendieck comparison defines $Q$. The same diagram construction for concatenated ordinals, with any finite number of marked junctions, supplies composition compatibility and its coherences.

On restricted hammocks this construction is exactly inversion of backward arrows and coherent composition, so $QBJ$ is the canonical localization comparison represented by $\rho$. The universal localization property, with its structural functor from $M$ retained, fixes this comparison; no arbitrary equivalence of mapping spaces is substituted. B.2 shows that $QBJ$ is a weak equivalence. Inverting it in $\mathrm{Hot}$ and composing with $Q$ gives a left inverse of $BJ$. The assertion about homotopy groups follows, with the path from the chosen compatible comparison providing the basepoint identification.

For completeness, the component assertion does not require higher comparisons. The effect of a path is its word in the ordinary localization. Identity insertion, same-direction composition, and cancellation of a marked arrow against its reversed traversal are realized by path morphisms or zigzags of them. A commuting square at a backward arrow uses only inversion of that backward arrow; no vertical component needs to be inverted. These are precisely the relations in ordinary localization. Thus the effect is surjective and identifies exactly the zigzag components of the full category. The restricted assertion is Dwyer--Kan Proposition~3.1. The component identifications commute with $J$.
\end{proof}

\begin{result}{B.7}{Proved}{The elementary grid move and its limited consequence}
A commuting square with backward horizontal arrows in $\Sigma$ and arbitrary vertical arrows gives, after identity padding, a restricted hammock between the two mixed boundary paths. This does not require two-out-of-three.
\end{result}
\begin{proof}
Write the square as
\[
 w:x'\to x,\qquad w':y'\to y,\qquad
 v:x\to y,\qquad v':x'\to y',\qquad vw=w'v',
 \quad w,w'\in\Sigma.
\]
The two paths $x\xleftarrow{w}x'\xrightarrow{v'}y'$ and
$x\xrightarrow{v}y\xleftarrow{w'}y'$ can be padded to type $(-,+,-)$ as
\[
 x\xleftarrow{w}x'\xrightarrow{v'}y'\xleftarrow{1}y',
 \qquad
 x\xleftarrow{1}x\xrightarrow{v}y\xleftarrow{w'}y'.
\]
The component list $(1_x,w,w',1_{y'})$ is a restricted natural transformation from the first padded path to the second. Its middle square is the given equation; its other squares are identity equations. For a forward horizontal square, the two mixed routes compose to the same arrow and are related by type contractions with identity components. Sweeping these local moves through a finite ladder proves its equality on components using restricted paths. A proof of a weak equivalence of entire nerves would additionally have to organize these sweeps on all simplices, compatibly with faces, degeneracies, and compositions. The two-dimensional square calculation itself does not supply that missing assertion.
\end{proof}

\begin{result}{B.8}{Open}{The unrestricted-vertical inclusion beyond the proved cases}
This paper does not settle whether the specified functor $J$ is in $W_\infty$ for every wide $\Sigma$, for every $\Sigma$ satisfying two-out-of-three, or for every saturated $\Sigma$. It also does not settle that assertion for all relative categories admitting a homotopy calculus of fractions, all model categories, or all Brown categories of fibrant objects. The special cases proved in B.4 and the exact short-fraction statements in B.5 retain their stated scope. No counterexample to the general weak-equivalence assertion is established here.
\end{result}
\begin{proof}[Status and exact remaining step]
B.2 identifies the restricted side for every pair. B.6 gives a canonical retraction on homotopy types and the correct components on the full side; it does not prove that the full side has no additional higher homotopy. B.7 removes each arbitrary ladder at the level of one-dimensional homotopy classes. The attempted general argument would require a contractible category of compatible grid sweeps for every higher simplex and compatibility under all simplex operations. That coherent replacement construction has not been proved here. In the model and Brown cases, reducing the restricted mapping space to a short fraction category leaves the same additional comparison into the full all-type category; a theorem about the restricted category cannot be used to close that gap.

A concrete proper marked example of the smallest object size remaining in our tested family, not covered by B.4 and left undecided by these arguments, is the chain $0<1<2$ with identities and $0\to2$ marked. It satisfies two-out-of-three. Its marked arrows are saturated: the functor to the ordinary category generated by a split inclusion $i:A\to B$ and retraction $r:B\to A$ with $ri=1_A$ and $ir\ne1_B$, sending $0,2$ to $A$, $1$ to $B$, and the two generating arrows to $i,r$, inverts the marked composite while neither generator becomes invertible. Thus no additional arrow of $[2]$ is forced invertible. This is an undecided instance in the present proof, not a proposed counterexample. Our finite computations in B.9 give no discrepancy for it.

The literature search included Dwyer--Kan's \emph{Calculating simplicial localizations} and \emph{Function complexes in homotopical algebra}, Arakawa--Cnossen's manuscript, and searches for ``hammock arbitrary vertical maps'' and ``Grothendieck path categories simplicial localization''. The search located the same question in a 2020 discussion, but no checked primary theorem closing the unrestricted-vertical step. We make no assertion that it is an open problem in the published literature, and no claim of priority.
\end{proof}

\begin{result}{B.9}{Proved}{What the finite search establishes}
For the explicitly specified search family and bounds in
\texttt{code/b\_path\_search.py}, the full and restricted path-category truncations have equal $H_0(-;\mathbb F_2)$ and $H_1(-;\mathbb F_2)$. There are 66 relative-category-and-endpoint cases at each of word lengths three and four, for 132 comparisons in total. This finite result proves no unbounded comparison theorem and supplies no counterexample to B.8.
\end{result}
\begin{proof}
The program considers the two-element chain, three-element chain, span, cospan, and four-element crown. In each it enumerates all composition-closed wide markings and the two endpoint pairs $(0,0)$ and $(0,n-1)$, and records whether two-out-of-three holds. For every word of length at most the selected bound it enumerates all relative paths. A morphism retains its full endpoint-preserving type map; two distinct type maps are never identified merely because their source and target paths agree. Poset naturality is tested componentwise, and composites are calculated by composing the vertex maps. Thus all arrows and all composable pairs of the finite truncation are included.

To compute $H_1$ over $\mathbb F_2$, choose a spanning forest of its directed one-skeleton viewed as an undirected multigraph. The remaining edges form generators after the forest is collapsed. Every composable pair contributes the relation consisting of its two edges and its composite, with identities omitted. Gaussian elimination on these binary relations computes the quotient dimension. These are exactly the normalized nerve chain relations in degrees one and two. The code separately records the number of objects, nonidentity arrows, nondegenerate two-simplices, components, and the resulting first Betti number. The two complete output files are \texttt{b\_path\_search\_L3.jsonl} and \texttt{b\_path\_search\_L4.jsonl}; all 132 comparisons agree. Higher-length paths can create or kill homology, and the outputs make no persistence claim.
\end{proof}
